\documentclass{article}
\usepackage{graphicx}
\usepackage{authblk}
\usepackage{indentfirst}
\usepackage{amssymb}
\usepackage{amsthm}
\usepackage{amsmath}
\usepackage{hyperref}
\usepackage{enumitem}
\usepackage{pgfplots}
\usepackage{tikz}
\usetikzlibrary{arrows.meta,calc}
\pgfplotsset{compat=1.18}
\newtheorem{theorem}{Theorem}
\newtheorem{corollary}{Corollary}[theorem]
\newtheorem{lemma}[theorem]{Lemma}
\newtheorem{definition}{Definition}
\newtheorem{problem}{Problem}
\newtheorem{solution}{Solution}
\newtheorem*{example}{Example}
\newtheorem{remark}{Remark}
\newtheorem{proposition}{Proposition}
\newtheorem{method}{Method}
\newtheorem{principle}{Principle}
\newcommand{\Span}{\operatorname{span}}
\newenvironment{amatrix}[1]{%
  \left(\begin{array}{@{}*{#1}{c}|c@{}}
}{%
  \end{array}\right)
}
\reversemarginpar
\hypersetup{colorlinks=true,}
\begin{document}

%------- Title page -----------
\title{MATH 323: Probability}
\author{W.R.YCH.}
\date{May $4^{th}$, 2026}
\setcounter{Maxaffil}{0}
\renewcommand\Affilfont{\itshape\small}
\maketitle

%------- Abstract -----------
\noindent\rule{\textwidth}{0.4pt}
\thispagestyle{empty}
\begin{abstract}
\end{abstract}
\noindent\rule{\textwidth}{0.4pt}
\clearpage

%------- Table of Contents -----------
\thispagestyle{empty}
{
  \hypersetup{linkcolor=black}
  \tableofcontents
}
\clearpage

%------- Introduction -----------
\setcounter{page}{1}
\section{Introduction}
These notes were written for MATH 323 (Probability) at McGill University, Summer 2026. The course is taught by Magid Sabbagh and meets in Stewart Biology Building S1/3 every Monday, Tuesday, Wednesday, and Thursday from 8:35\,am to 10:55\,am.

\subsubsection*{Grading scheme}
The final grade is the higher of the two:
\begin{align}
    &20\%\ \text{assignments} + 5\%\ \text{multiple choice quiz} + 25\%\ \text{midterm} + 50\%\ \text{final}\\
    &\text{or}\\
    &20\%\ \text{assignments} + 5\%\ \text{multiple choice quiz} + 75\%\ \text{final}
\end{align}
The textbook is \textit{Mathematical Statistics with Applications}, 7th edition. The midterm is 1h30 long but students get the full 2h30 class block to complete it.

% ============================================================
\section{Prerequisite Knowledge}
% ============================================================
\marginpar{May 4, 2026}

\subsection{Sets and Set Operations}

A \textbf{set} is a well-defined collection of objects, where ``well-defined'' means we can determine whether or not a given element $x$ belongs to the set. We write $a \in A$ if $a$ is an element of $A$, and $a \notin A$ otherwise. A set is specified either by listing its elements or by giving a property that defines membership.

We always work inside a fixed \textbf{universal set} $\Omega$. Let $A, B \subseteq \Omega$. We define the following operations:
\begin{itemize}
    \item \textbf{Subset}: $A \subseteq B$ if $x \in A$ implies $x \in B$.
    \item \textbf{Equality}: $A = B$ if $A \subseteq B$ and $B \subseteq A$.
    \item \textbf{Proper subset}: $A \subsetneq B$ if $A \subseteq B$ and $A \neq B$.
    \item \textbf{Union}: $A \cup B = \{x : x \in A \text{ or } x \in B\}$.
    \item \textbf{Intersection}: $A \cap B = \{x : x \in A \text{ and } x \in B\}$.
    \item \textbf{Difference}: $A \setminus B = \{x : x \in A \text{ and } x \notin B\}$.
    \item \textbf{Complement}: $A^c = \{x : x \in \Omega \text{ and } x \notin A\}$, also written $\overline{A}$.
    \item \textbf{Singleton}: $\{x\} = \{y : y = x\}$.
\end{itemize}
The empty set is denoted $\emptyset$.

\begin{proposition}
For any $A \subseteq \Omega$:
\[
A \cup A^c = \Omega, \qquad A \cap A^c = \emptyset, \qquad A \cup \emptyset = A, \qquad A \cap \emptyset = \emptyset.
\]
\end{proposition}

\begin{proposition}
Let $A, B, C \subseteq \Omega$. Then:
\begin{enumerate}
    \item \textbf{Commutativity}: $A \cup B = B \cup A$ and $A \cap B = B \cap A$.
    \item \textbf{Associativity}: $A \cup (B \cup C) = (A \cup B) \cup C$ and $A \cap (B \cap C) = (A \cap B) \cap C$.
    \item \textbf{Distributivity}: $A \cup (B \cap C) = (A \cup B) \cap (A \cup C)$ and $A \cap (B \cup C) = (A \cap B) \cup (A \cap C)$.
\end{enumerate}
\end{proposition}

\begin{proof}
We prove $A \cap (B \cup C) = (A \cap B) \cup (A \cap C)$; the other identities are similar.

\textbf{($\subseteq$)} Let $x \in A \cap (B \cup C)$. Then $x \in A$ and $x \in B \cup C$, so $x \in B$ or $x \in C$. If $x \in B$, then $x \in A \cap B$. If $x \in C$, then $x \in A \cap C$. In either case $x \in (A \cap B) \cup (A \cap C)$.

\textbf{($\supseteq$)} Let $x \in (A \cap B) \cup (A \cap C)$. If $x \in A \cap B$, then $x \in A$ and $x \in B \subseteq B \cup C$, so $x \in A \cap (B \cup C)$. The case $x \in A \cap C$ is analogous.
\end{proof}

\begin{proposition}[De Morgan's Laws]
For $A, B \subseteq \Omega$:
\[
(A \cup B)^c = A^c \cap B^c, \qquad (A \cap B)^c = A^c \cup B^c.
\]
\end{proposition}

\subsection{Indexed Unions and Intersections}

Let $A_1, A_2, A_3, \ldots$ be subsets of $\Omega$. For $n \in \mathbb{N}$:
\begin{align*}
\bigcap_{i=1}^n A_i &= \{x \in \Omega : x \in A_i \text{ for all } i \in \{1, \ldots, n\}\},\\
\bigcup_{i=1}^n A_i &= \{x \in \Omega : x \in A_i \text{ for at least one } i \in \{1, \ldots, n\}\}.
\end{align*}
The countable versions are defined the same way with $i$ ranging over $\mathbb{N}$:
\[
\bigcap_{i=1}^\infty A_i = \{x \in \Omega : x \in A_i \text{ for all } i \in \mathbb{N}\}, \qquad \bigcup_{i=1}^\infty A_i = \{x \in \Omega : x \in A_i \text{ for some } i \in \mathbb{N}\}.
\]

\begin{proposition}[Generalized De Morgan]
\[
\left(\bigcup_{i=1}^\infty A_i\right)^c = \bigcap_{i=1}^\infty A_i^c, \qquad \left(\bigcap_{i=1}^\infty A_i\right)^c = \bigcup_{i=1}^\infty A_i^c.
\]
\end{proposition}

\subsection{Cardinality and Inclusion-Exclusion}

\begin{problem}
In a class of 101 students, 20 play hockey, 30 like math, and 5 do both. How many students:
\begin{enumerate}
    \item do not play hockey?
    \item do not like math?
    \item neither play hockey nor like math?
    \item play hockey but do not like math?
\end{enumerate}
\end{problem}

\begin{solution}
\begin{enumerate}
    \item $101 - 20 = 81$.
    \item $101 - 30 = 71$.
    \item $101 - (20 + 30 - 5) = 56$.
    \item $20 - 5 = 15$.
\end{enumerate}
\end{solution}

We have the identity $A \setminus B = A \cap B^c = A \setminus (A \cap B)$.

\begin{remark}
Write $|A|$ for the cardinality of $A$. Then
\[
|A \cup B| = |A| + |B| - |A \cap B|.
\]
\end{remark}

\begin{proof}
Partition $A \cup B$ into three disjoint pieces: $A \setminus B$, $A \cap B$, and $B \setminus A$. Let $x = |A \setminus B|$, $y = |A \cap B|$, and $z = |B \setminus A|$. Then $|A \cup B| = x + y + z$, $|A| = x + y$, and $|B| = y + z$, so
\[
|A| + |B| - |A \cap B| = (x + y) + (y + z) - y = x + y + z = |A \cup B|. \qedhere
\]
\end{proof}

% ============================================================
\section{Sigma-Algebras and Probability Measures}
% ============================================================
\marginpar{May 5, 2026}

\subsection{Sigma-Algebras}

\begin{definition}
A \textbf{$\sigma$-algebra} on a set $\Omega$ is a collection $\mathbb{F}$ of subsets of $\Omega$ such that:
\begin{enumerate}
    \item $\emptyset \in \mathbb{F}$,
    \item if $A \in \mathbb{F}$, then $A^c \in \mathbb{F}$,
    \item if $A_1, A_2, \ldots \in \mathbb{F}$, then $\bigcup_{i=1}^\infty A_i \in \mathbb{F}$.
\end{enumerate}
Elements of $\mathbb{F}$ are called \textbf{events}.
\end{definition}

\begin{example}\
\begin{enumerate}
    \item Let $\Omega$ be any set. Then $\mathbb{F} = \{\emptyset, \Omega\}$ is a $\sigma$-algebra (the trivial one).
    \item Let $\Omega$ be any set. The power set $2^\Omega$ (the collection of all subsets of $\Omega$) is a $\sigma$-algebra.
    \item Let $\Omega$ be a set. Then $\mathbb{F} = \{A \subseteq \Omega : A \text{ is countable or } A^c \text{ is countable}\}$ is a $\sigma$-algebra.
    \item For any $A \subseteq \Omega$, the collection $\mathbb{F} = \{\emptyset, A, A^c, \Omega\}$ is a $\sigma$-algebra.
\end{enumerate}
\end{example}

\begin{example}
Flip a coin once. Then $\Omega = \{H, T\}$ and $2^\Omega = \{\emptyset, \{H\}, \{T\}, \{H, T\}\}$.
\end{example}

\subsection{Probability Measures}

\begin{definition}
Let $\Omega$ be the sample space of an experiment and $\mathbb{F}$ a $\sigma$-algebra of subsets of $\Omega$. A \textbf{probability measure} is a function $P : \mathbb{F} \to [0, 1]$ satisfying the \textbf{Kolmogorov axioms}:
\begin{enumerate}
    \item $P(A) \geq 0$ for all $A \in \mathbb{F}$.
    \item $P(\Omega) = 1$.
    \item \textbf{Countable additivity}: if $A_1, A_2, \ldots \in \mathbb{F}$ are pairwise disjoint (that is, $A_k \cap A_l = \emptyset$ whenever $k \neq l$), then
    \[
    P\!\left(\bigcup_{i=1}^\infty A_i\right) = \sum_{i=1}^\infty P(A_i).
    \]
\end{enumerate}
The triple $(\Omega, \mathbb{F}, P)$ is called a \textbf{probability space}.
\end{definition}

\subsection{Basic Properties of Probability}

\begin{proposition}[Monotonicity]
If $A \subseteq B$, then $P(A) \leq P(B)$.
\end{proposition}

\begin{proof}
Write $B = A \cup (B \setminus A)$ as a disjoint union. By countable additivity, $P(B) = P(A) + P(B \setminus A) \geq P(A)$, since $P(B \setminus A) \geq 0$.
\end{proof}

\begin{corollary}
For any event $A$, $P(A) \leq 1$.
\end{corollary}

\begin{proof}
Since $A \subseteq \Omega$, monotonicity gives $P(A) \leq P(\Omega) = 1$.
\end{proof}

\begin{proposition}[Complement Rule]
$P(A^c) = 1 - P(A)$.
\end{proposition}

\begin{proof}
Since $\Omega = A \cup A^c$ disjointly, $1 = P(\Omega) = P(A) + P(A^c)$, so $P(A^c) = 1 - P(A)$.
\end{proof}

\begin{proposition}[Inclusion-Exclusion for Two Events]
For any events $A, B \in \mathbb{F}$:
\[
P(A \cup B) = P(A) + P(B) - P(A \cap B).
\]
\end{proposition}

\begin{proof}
Decompose $A \cup B = A \cup (B \cap A^c)$ where $A$ and $B \cap A^c$ are disjoint. By countable additivity,
\[
P(A \cup B) = P(A) + P(B \cap A^c). \tag{$*$}
\]
Similarly, $B = (A \cap B) \cup (A^c \cap B)$ is a disjoint decomposition, so
\[
P(B) = P(A \cap B) + P(A^c \cap B), \qquad \text{i.e.}\qquad P(A^c \cap B) = P(B) - P(A \cap B).
\]
Substituting this into $(*)$ gives $P(A \cup B) = P(A) + P(B) - P(A \cap B)$.
\end{proof}

\subsection{Boole's Inequality and Bonferroni}

The countable additivity axiom requires the events to be pairwise disjoint. What can we say about $P(\bigcup_{i=1}^\infty A_i)$ when the $A_i$ are not disjoint?

\begin{proposition}[Boole's Inequality / Union Bound]
For any sequence of events $\{A_i\}_{i=1}^\infty$,
\[
P\!\left(\bigcup_{i=1}^\infty A_i\right) \leq \sum_{i=1}^\infty P(A_i).
\]
\end{proposition}

\begin{proof}
Define a new sequence by ``carving out'' the previously seen sets:
\[
A_1^\star = A_1, \qquad A_n^\star = A_n \setminus \bigcup_{i=1}^{n-1} A_i \quad \text{for } n \geq 2.
\]

\textbf{Step 1: The $A_i^\star$ are pairwise disjoint.} Suppose for contradiction that $A_k^\star \cap A_l^\star \neq \emptyset$ with $k < l$. Pick $x$ in the intersection. Since $x \in A_l^\star$, by definition $x \notin \bigcup_{i=1}^{l-1} A_i$, in particular $x \notin A_k$. But $x \in A_k^\star \subseteq A_k$, contradiction.

\textbf{Step 2: $\bigcup_{i=1}^\infty A_i^\star = \bigcup_{i=1}^\infty A_i$.} Since $A_i^\star \subseteq A_i$ for all $i$, the inclusion $\bigcup A_i^\star \subseteq \bigcup A_i$ is immediate. For the reverse, let $x \in \bigcup A_i$ and let $i_0$ be the smallest index with $x \in A_{i_0}$. Then $x \notin A_i$ for $i < i_0$, so $x \in A_{i_0}^\star \subseteq \bigcup A_i^\star$.

\textbf{Step 3: Conclude.} By countable additivity applied to the $A_i^\star$ together with monotonicity ($A_i^\star \subseteq A_i$ implies $P(A_i^\star) \leq P(A_i)$),
\[
P\!\left(\bigcup_{i=1}^\infty A_i\right) = P\!\left(\bigcup_{i=1}^\infty A_i^\star\right) = \sum_{i=1}^\infty P(A_i^\star) \leq \sum_{i=1}^\infty P(A_i). \qedhere
\]
\end{proof}

\begin{corollary}[Bonferroni Inequality]
For events $A_1, \ldots, A_n$,
\[
P\!\left(\bigcap_{i=1}^n A_i\right) \geq \sum_{i=1}^n P(A_i) - (n-1).
\]
\end{corollary}

\begin{proof}
Using De Morgan's law, the complement rule, and the union bound,
\begin{align*}
P\!\left(\bigcap_{i=1}^n A_i\right) &= 1 - P\!\left(\bigcup_{i=1}^n A_i^c\right)\\
&\geq 1 - \sum_{i=1}^n P(A_i^c)\\
&= 1 - \sum_{i=1}^n \bigl(1 - P(A_i)\bigr)\\
&= \sum_{i=1}^n P(A_i) - (n-1). \qedhere
\end{align*}
\end{proof}

\begin{problem}
Suppose $\{A_i\}_{i=1}^\infty$ are events with $P(A_i) = 1$ for all $i \in \mathbb{N}$. Show that $P\!\left(\bigcap_{i=1}^\infty A_i\right) = 1$.\footnote{He'll prolly ask this on an exam.}
\end{problem}

\subsection{Independent Events}

\begin{definition}
Two events $A$ and $B$ are \textbf{independent} if $P(A \cap B) = P(A) \cdot P(B)$.
\end{definition}

% ============================================================
\section{Finite Sample Spaces and Equally Likely Outcomes}
% ============================================================

\begin{example}
Roll a die twice, with the rolls independent, and record the pair of outcomes. The sample space is
\[
\Omega = \{(a, b) : a, b \in \{1, 2, 3, 4, 5, 6\}\}, \qquad |\Omega| = 36.
\]
By independence, $P(\{(a, b)\}) = \tfrac{1}{6} \cdot \tfrac{1}{6} = \tfrac{1}{36}$, so all 36 outcomes are equally likely.
\end{example}

\begin{definition}
Let $\Omega = \{w_1, \ldots, w_N\}$ be a finite sample space with $|\Omega| = N$. The outcomes are \textbf{equally likely} if
\[
P(\{w_1\}) = P(\{w_2\}) = \cdots = P(\{w_N\}) = \tfrac{1}{N}.
\]
\end{definition}

For any event $A \subseteq \Omega$ with $|A| = d$, writing $A = \{w_{i_1}, \ldots, w_{i_d}\}$,
\[
P(A) = \sum_{k=1}^d P(\{w_{i_k}\}) = \frac{d}{N} = \frac{|A|}{|\Omega|}.
\]

\begin{example}
A fair coin is flipped twice independently. The sample space is $\Omega = \{HH, HT, TH, TT\}$ with each outcome having probability $\tfrac{1}{4}$.

\textit{What is the probability of getting at least one head?} Let $A$ be that event. Then $A = \{HH, HT, TH\}$, so $|A| = 3$ and $P(A) = \tfrac{3}{4}$.
\end{example}

% ============================================================
\section{Counting}
% ============================================================

\begin{theorem}[Fundamental Counting Principle]
Given $m$ elements $a_1, \ldots, a_m$ and $n$ elements $b_1, \ldots, b_n$, the number of distinct ordered pairs $(a, b)$ with $a \in \{a_1, \ldots, a_m\}$ and $b \in \{b_1, \ldots, b_n\}$ is $mn$.
\end{theorem}

\begin{example}
How many 7-digit telephone numbers can be formed?

Each digit has 10 choices, giving $10^7$ numbers.

\textit{How many if the third digit must be even?} The third digit has 5 choices ($0, 2, 4, 6, 8$) and each of the other six digits has 10 choices, giving $10^6 \cdot 5$.
\end{example}

\subsection{Permutations}

\begin{example}
In a class of 100 students, 2 are chosen: the first wins \$100 and the second wins \$75. How many choices are there if:
\begin{enumerate}
    \item a student can win twice (\textbf{sampling with replacement})?
    \item a student can win at most once (\textbf{sampling without replacement})?
\end{enumerate}
\end{example}

\begin{solution}
\begin{enumerate}
    \item 100 choices for the first prize and 100 for the second: $100^2 = 10{,}000$.
    \item 100 choices for the first prize and 99 for the second: $100 \cdot 99 = 9{,}900$.
\end{enumerate}
\end{solution}

\begin{definition}
An ordered arrangement of $r$ distinct objects chosen from $n$ objects (without replacement) is called a \textbf{permutation}. The number of such permutations is
\[
P^n_r = n(n-1)(n-2) \cdots (n - r + 1) = \frac{n!}{(n-r)!}.
\]
\end{definition}

\begin{example}
In a class of 100 students, 2 students are chosen to form a committee where order does not matter. How many committees are possible?

There are $100 \cdot 99$ ordered pairs, but each unordered committee is counted $2!$ times, so the number of committees is
\[
\frac{100 \cdot 99}{2!} = 4{,}950.
\]
\end{example}

% ============================================================
\section{Appendix}
% ============================================================


% ============================================================
\section{Solutions}
% ============================================================


% ============================================================
\section{Useful Links}
% ============================================================

\end{document}